\documentclass[11pt,a4paper]{article}
\usepackage[top=15mm,bottom=15mm,left=16mm,right=16mm]{geometry}
\usepackage[T1]{fontenc}
\usepackage[utf8]{inputenc}
\usepackage{times}
\usepackage{xcolor}
\definecolor{DeepBlue}{HTML}{003366}
\definecolor{AccentBlue}{HTML}{0B5C8C}
\definecolor{Gold}{HTML}{B8860B}
\definecolor{LightBlue}{HTML}{EAF2F8}
\definecolor{LightGold}{HTML}{FBF3E0}
\usepackage[colorlinks=true,urlcolor=AccentBlue,linkcolor=DeepBlue,citecolor=DeepBlue]{hyperref}
\usepackage{titlesec}
\usepackage{enumitem}
\usepackage{amsmath,amssymb}
\usepackage{booktabs}
\usepackage{tabularx}
\usepackage{array}
\usepackage{colortbl}
\usepackage{multirow}
\usepackage[activate={true,nocompatibility},protrusion=true,expansion=false]{microtype}

\titleformat{\section}{\normalsize\bfseries\color{DeepBlue}}{\thesection.}{0.4em}{}
[\vspace{-0.35em}{\color{Gold}\rule{\linewidth}{0.6pt}}\vspace{0.1em}]
\titleformat{\subsection}{\small\bfseries\color{AccentBlue}}{\thesubsection}{0.4em}{}
\titleformat{\subsubsection}[runin]{\small\bfseries}{}{0em}{}[.]
\titlespacing*{\section}{0pt}{0.7em}{0.25em}
\titlespacing*{\subsection}{0pt}{0.45em}{0.12em}

\setlength{\parindent}{0pt}
\setlength{\parskip}{0.35em}
\pagestyle{plain}
\setlist{leftmargin=1.1em,itemsep=0.1em,topsep=0.12em,parsep=0em}
\newcolumntype{Y}{>{\raggedright\arraybackslash}X}
\renewcommand{\arraystretch}{1.15}

\begin{document}


\begin{center}
{\large\bfseries\color{DeepBlue} CERTIFLIGHT}\\[0.15em]
{\bfseries\color{AccentBlue} Certified End-to-End Autonomy Assurance for Civil UAV Missions:\\
From Network Detection through Navigation Certificates to Constrained Swarm Control}\\[0.35em]
{\footnotesize MSCA Postdoctoral Fellowship (European Fellowship) --- Part B1 working draft\\
\textbf{Fellow:} Roger Nick Anaedevha \quad$\cdot$\quad \textbf{Host:} SPRITZ Group, University of Padua \quad$\cdot$\quad
\textbf{Supervisor:} Prof.\ Mauro Conti \quad$\cdot$\quad \textbf{Co-supervisor:} Dr.\ Federico Cor\`o\\
\textbf{Duration:} 24 months \quad$\cdot$\quad \textbf{Panel:} ENG / PE6 \quad$\cdot$\quad \textbf{Call deadline:} 9 September 2026}
\end{center}

\vspace{-0.2em}
{\footnotesize\itshape Draft for review by the host. Structured to the MSCA-PF Part~B1 template (Excellence 50\,\%, Impact 30\,\%, Implementation 20\,\%; 10-page limit). Host-supplied text marked [Host]. This revision incorporates the co-supervisor's guidance: the swarm work package is now formulated as \emph{constrained reinforcement learning} rather than a combinatorial covering game; the threat model is broadened to non-malicious disturbances; and the impact case is built on civil, non-military use cases.}

% ==========================================================
\section{Excellence}

\subsection{Quality and pertinence of the research objectives; ambition and advance beyond the state of the art}

\subsubsection{The problem}
Civil autonomous UAVs are becoming infrastructure. Automated external defibrillators are already delivered by drone ahead of ambulances in Swedish emergency medical services; blood and diagnostic samples move by drone between UK and Scottish hospitals; drones monitor air quality, detect wildfire ignition, survey floods, and inspect energy and rail networks. Europe's commercial drone market was valued in the multi-billion-euro range in 2024 and is projected to roughly double by 2030. Every one of these missions depends on a learned autonomy stack that must keep flying when its inputs degrade.

Those inputs degrade routinely, and from three distinct causes that are today handled by three disconnected communities:
\begin{enumerate}
\item \textbf{Network-layer attack.} The multi-hop swarm mesh, the command-and-control link and the on-board bus admit time-delay forwarding, false-data injection and denial attacks. Network intrusion detection catches much of this --- but a detector must tolerate natural jitter, so an adversary that stays just under the detector's tolerance is invisible \emph{by construction}.
\item \textbf{Sensing-layer attack.} GNSS denial and spoofing degrade the position solution. This is now a documented \emph{civil} hazard, not a military one: tens of thousands of satellite-navigation disruption reports have been logged by civil aircraft over the Baltic region since 2023, prompting European aviation-safety information bulletins, with position errors of tens of metres and multi-hour outages.
\item \textbf{Non-malicious disturbance.} Sensor degradation and IMU drift, partial actuator or rotor failure, and severe environmental forcing --- gusts, turbulence, obstacle wake --- perturb the same control loop. The fault-tolerant-control literature treats these as bounded disturbances with unknown bound; the security literature ignores them.
\end{enumerate}

Each community bounds its own layer and leaves the interface to assumption. Certified robustness bounds what a model does with a perturbation but is silent on where the perturbation comes from or how often it arrives. Intrusion detection bounds what reaches the aircraft but says nothing about the flight. Fault-tolerant control rejects disturbances but is not composed with either. \textbf{The property an operator, an insurer and a regulator actually need --- \emph{will this mission complete, and can we prove it?} --- falls in the gaps between them.}

\subsubsection{The idea}
CERTIFLIGHT unifies the three causes into a \textbf{single perturbation budget} and then \emph{composes} guarantees across the layer boundary. The central object is a chained statement
\[
\theta \;\longmapsto\; \Delta(\theta) \;\longmapsto\; \delta(\theta) \;\longmapsto\; \mathrm{MCR} \;\ge\; f\!\left(\delta_{\text{tot}}\right),
\qquad
\delta_{\text{tot}} \;=\; \delta_{\text{net}}(\theta) \;\oplus\; \delta_{\text{gnss}} \;\oplus\; \delta_{\text{fault}} \;\oplus\; \delta_{\text{env}},
\]
read: a network detector operating at point $\theta$ bounds the residual undetected attack $\Delta(\theta)$; that residual maps to a bounded perturbation $\delta(\theta)$ on the navigation input; the malicious, fault and environmental components combine into a total budget $\delta_{\text{tot}}$; and a certificate turns the total budget into a guaranteed \emph{Mission-Completion-Rate} (MCR) floor. Because the budget is source-agnostic, the same certificate covers an evading attacker, a drifting IMU and a gust --- which is precisely what an airworthiness argument requires.

The programme then asks the question that a single-aircraft certificate cannot answer: \textbf{what does a fleet do when an individual aircraft's certificate degrades?} Under jamming a UAV's certified navigation guarantee shrinks; below a safety threshold that aircraft can no longer be trusted with its assigned task. CERTIFLIGHT formulates swarm task allocation as a \textbf{Constrained Markov Decision Process} in which the per-agent certified bound is a \emph{hard constraint}, and learns a policy that dynamically re-allocates the fleet --- rerouting healthy aircraft to preserve coverage --- while provably respecting those constraints during both learning and deployment.

\subsubsection{Objectives}
\begin{itemize}
\item[\textbf{O1}] \textbf{Unified perturbation budget and composed certificate.} Prove a composition theorem chaining a network detector's operating point to a certified MCR floor, over a budget that aggregates malicious, fault-induced and environmental perturbations; calibrate the interface empirically on real flight data. (WP2)
\item[\textbf{O2}] \textbf{Open cross-layer benchmark and evaluation platform.} Release the first reproducible benchmark whose unit of analysis is the \emph{pairing} of a network or environmental event with the navigation degradation it induces, with auditable provenance. (WP3)
\item[\textbf{O3}] \textbf{Certified-constrained multi-agent reinforcement learning for swarm re-allocation.} Develop a constrained MARL formulation in which per-agent certified navigation radii are hard constraints, so that when jamming degrades an agent's guarantee the policy re-allocates the fleet to maintain mission coverage with feasibility guarantees. (WP4)
\item[\textbf{O4}] \textbf{Assurance evidence for civil certification.} Map the resulting certificates onto EASA SORA / U-space assurance levels, EU AI Act Article~15 (accuracy, robustness, cybersecurity), and DO-326A\,/\,ED-202A airworthiness security, so that a certified floor becomes usable compliance evidence. (WP5)
\end{itemize}

\subsubsection{Advance beyond the state of the art}
\emph{(i) Certified robustness} provides deterministic Lipschitz-type bounds and probabilistic randomized-smoothing radii for models, and has been lifted to policies --- CROP certifies a cumulative-reward lower bound for RL under bounded observation perturbation, CAROL extends certification to model-based abstract interpretation. All of this is \textbf{single-agent and observation-level}: it certifies a policy against a perturbation ball, never against a network operating point, a hardware fault, or a fleet-level coverage requirement.
\emph{(ii) Safe reinforcement learning} offers Constrained Policy Optimization and Lagrangian variants (TRPO-/PPO-Lagrangian, RCPO) with zero-duality-gap results, and pointwise-safe mechanisms --- control barrier functions, shielding, verification-guided shielding. Crucially, CPO and Lagrangian methods enforce \emph{expected} constraint satisfaction over a horizon; they cannot supply the pointwise guarantee a safety-critical flight needs. Barrier and shielding methods do, but their constraint sets are hand-specified, not derived from a formal robustness certificate.
\emph{(iii) Safe multi-agent RL} for UAV swarms exists (decentralized multi-CBF safe MADDPG; hierarchical MARL with barrier functions; QMIX-family coordination), and RL under GNSS denial is emerging (RL formation control in GPS-denied flight; EKF-plus-DRL trajectory optimisation under jamming). None of it consumes a \emph{certified} per-agent bound.

\textbf{The gap, stated precisely.} No published work composes formally certified per-agent navigation guarantees into the hard-constraint set of a constrained multi-agent policy that re-allocates a swarm when agents degrade --- and none derives those per-agent guarantees from a budget that unifies adversarial, fault and environmental perturbation. CERTIFLIGHT closes both. Its three technical firsts are: the first formal chain from a network-security operating point to a navigation-safety floor (O1); the first benchmark whose records are cross-layer pairings with auditable provenance (O2); and the first constrained-MARL swarm controller whose constraints are certified robustness bounds rather than hand-tuned safety margins (O3).

\subsection{Soundness of the methodology}

\subsubsection{WP2 --- Unified budget and composed certificates}
\emph{Formal core.} Let navigation state evolve as $\dot x = F(x,u)$ with $F$ $L$-Lipschitz over the operating region and mission horizon $T$. A bounded input perturbation $\lVert\delta u\rVert \le \delta_{\text{tot}}$ yields, via the integral form of Gr\"onwall's inequality, a trajectory-tube radius $\rho = \delta_{\text{tot}}e^{LT}$; the mission completes whenever $\rho$ fits inside the mission's safe-corridor margin $m$, giving the certified floor $f(\delta_{\text{tot}}) = \Pr[\delta_{\text{tot}}e^{LT}\le m]$. \emph{Interface.} On the network side, a detector tolerating per-hop residual $\theta$ admits cumulative undetected staleness capped by the store-and-forward contact-window slack; a measured delay-to-position rate converts staleness into position error, and this measured rate is an order of magnitude tighter than the kinematic worst case --- the difference between a vacuous and a useful certified window. \emph{Unification.} Fault and environmental components enter through the same channel: actuator and sensor faults and gust forcing are modelled, following the fault-tolerant-control literature, as bounded additive disturbances on the same input, and are aggregated with the malicious component by a union bound over independently certified sources. \emph{Coverage.} Because the deterministic bound degrades as $e^{LT}$, three complementary certificates cover the regimes it cannot: randomized smoothing where the budget is known only statistically, PAC-Bayes for the train-to-deployment distribution gap, and multiplicative-weights regret against an adaptive jammer. \emph{Validation.} Soundness is verified empirically --- the certified floor must lower-bound the measured completion rate at every operating point --- with multi-seed runs and Holm-corrected significance tests.

\subsubsection{WP3 --- Cross-layer benchmark and open platform}
A two-layer schema represents any observation from any source as a typed event, groups events into labelled attack or disturbance windows, and --- the novel object --- joins a network/environment window to the autonomy window it induces as a \emph{cross-layer pairing} carrying the detector operating point, residual budget, induced perturbation and mission outcome. Every pairing carries a mandatory provenance flag recording whether the coupling was measured on one platform, derived through a stated physical model, or aligned in simulation, so a reviewer can filter by evidence strength rather than trusting an unaudited join. The platform ships lossless ingestion adapters for established public corpora plus a physics-informed navigation benchmark, a detector-operating-curve harness that sweeps $\theta$ and emits recall/false-positive rate, and a reference pairing pipeline. Third-party data is fetched from origin under its own licence; the release contains our own data, adapters and unified labels.

\subsubsection{WP4 --- Certified-constrained multi-agent RL for swarm re-allocation}
\emph{Formulation.} The fleet is modelled as a Constrained Markov Decision Process $(\mathcal S,\mathcal A,P,r,\{c_i\},\{d_i\})$. Agent $i$'s state includes its \emph{current certified navigation radius} $R_i(t)$, computed online by the WP2 certificate from the locally estimated jamming and fault conditions. The task-allocation policy maximises mission coverage subject to per-agent constraints of the form $c_i:\;R_i(t)\ \ge\ R^{\min}(\tau_i)$, i.e.\ an aircraft may be assigned task $\tau_i$ only while its certified radius exceeds the radius that task provably requires. \emph{Solution strategy, in increasing strength of guarantee.} (a) A Lagrangian baseline (PPO-Lagrangian, CPO) gives expected-constraint satisfaction and a performance reference. (b) A \emph{certified safety layer} projects the policy's proposed allocation onto the feasible set defined by the current certified radii, giving pointwise constraint satisfaction at deployment --- the certificate replaces a hand-tuned barrier margin. (c) Shielding with a verified fallback allocation guarantees that when no feasible assignment exists, coverage degrades gracefully to a provably attainable subset rather than failing silently. \emph{Analysis targets.} Feasibility preservation under the safety layer; a bound on coverage loss as a function of the fraction of agents whose certified radius has collapsed; and the operational trade-off frontier between per-agent robustness investment and fleet size required for a target coverage guarantee. \emph{Evaluation.} Multi-agent simulation on the WP3 platform, sweeping jamming intensity and injected fault/gust profiles across civil mission geometries (urban corridor monitoring, environmental transects, point-to-point medical delivery), against jamming-blind and worst-case-margin baselines. \emph{Rationale for this formulation.} Constrained optimisation is the natural home for the project's identity --- provable bounds as hard constraints --- and it scales to the dynamic, partially observable conditions of real swarm flight in a way static combinatorial covering models do not.

\subsubsection{WP5 --- Assurance mapping (cross-cutting)}
Each certificate is expressed as evidence against three frameworks: EASA's Specific Operations Risk Assessment and the assurance levels it derives (the regime governing beyond-visual-line-of-sight civil operations); EU AI Act Article~15, which requires high-risk AI to achieve appropriate accuracy, robustness and cybersecurity, to be resilient to errors and faults, and resilient against attempts to alter its outputs --- language that maps almost directly onto a certified MCR floor over a unified perturbation budget; and DO-326A\,/\,ED-202A airworthiness security. Output: a white paper and a worked assurance case for one civil mission profile.

\subsubsection{Interdisciplinarity, open science, and the gender dimension}
The project deliberately spans machine-learning robustness, network security, control and constrained optimisation, and regulatory assurance; the fellow and host cover complementary halves and the co-supervisor brings the algorithms and UAV-networks expertise WP4 requires. \emph{Open science} is built in, not bolted on: pre-registration of experimental protocols where applicable, open-access publication (green route), release of the benchmark, adapters, harness and trained policies under permissive licences with provenance ledgers, and reproducible per-seed result files so that every figure can be regenerated. \emph{Gender and diversity:} the research object is a technical system, so no sex/gender dimension arises in the subject matter; the dimension is addressed in practice --- inclusive recruitment of any students supervised, and attention to equitable access in the civil use cases studied (e.g.\ emergency medical delivery to rural and under-served areas, where drone logistics has the largest marginal benefit).

\subsection{Quality of the supervision, training, and two-way knowledge transfer}
\textbf{Supervision.} Prof.\ Mauro Conti (supervisor) is among Europe's most cited security researchers and leads SPRITZ, a large group with an established record of hosting funded fellows. Dr.\ Federico Cor\`o (co-supervisor) leads the group's UAV algorithms and networks line and will co-supervise WP4 directly. Weekly one-to-one meetings with the co-supervisor, monthly progress reviews with the supervisor, and quarterly reviews against the Career Development Plan. [Host: formal supervision arrangements, mentoring, group seminar structure.]

\textbf{Training (to the fellow).} Scientific: constrained and multi-agent reinforcement learning, combinatorial and network algorithms, UAV network protocols and testbed practice --- competences the fellow does not currently hold at depth and which the host leads. Transferable: proposal writing and grant acquisition (an ERC Starting Grant concept is a project deliverable), student co-supervision, project and data management, science communication, and research ethics and security. [Host: doctoral-school courses, institutional training catalogue, language provision.]

\textbf{Transfer to the host (two-way).} The fellow brings certified-robustness machinery, a physics-informed navigation benchmark of over one hundred thousand simulated flights, four implemented certification methods, and production-grade open-source security platform engineering. These extend SPRITZ's UAV security line from empirical defence to formal guarantees --- a capability the group does not currently hold --- and are transferred through a group seminar series, co-supervised student projects, and shared code ownership of the released platform.

\subsection{Quality and appropriateness of the researcher's experience, competences and skills}
The fellow's doctoral work (NRNU MEPhI, defended June 2026, degree conferred December 2026) is on adversarially robust machine learning for intrusion detection and autonomous systems. Relevant published outputs include poisoning-resilient state-space models (\emph{Expert Systems with Applications}), a PAC-Bayesian adversarial-robustness framework (\emph{Knowledge and Information Systems}), Byzantine-resilient federated aggregation via stochastic games (\emph{Complex and Intelligent Systems}), and a first-place best-paper award at IEEE EDM 2026 for stochastic temporal graph networks for industrial intrusion detection (accepted, \emph{IEEE Transactions on Industry Applications}), with further manuscripts under review at IEEE Transactions. A joint cross-layer UAV-security manuscript with the host group is already in submission, evidencing an active and productive collaboration. Engineering capacity for WP3 is evidenced by deployed open-source security platforms and four years of prior production cybersecurity engineering, which also grounds the standards-facing orientation of WP5. The fellowship is the step from a strong doctoral record to independent investigator: it supplies the swarm-control and networks competences the fellow lacks, a European host and network, and the grant-writing experience required to lead a group.

% ==========================================================
\section{Impact}

\subsection{Credibility of the measures to enhance the fellow's career perspectives and employability}
The fellowship is designed around a single career objective: to establish the fellow as an independent European investigator at the autonomy-assurance interface, capable of leading a group. Four concrete measures. \emph{(i) Competence acquisition:} constrained and multi-agent RL and network algorithms (WP4, with the co-supervisor) close the principal gap in the fellow's profile, converting a certified-robustness specialist into a researcher who can address fleet-level autonomy. \emph{(ii) Independence:} the fellow leads all work packages, owns the released platform, and co-supervises at least one Master's and one doctoral student. \emph{(iii) Funding capability:} structured grant-writing training culminating in a submitted ERC Starting Grant concept and at least one national or European co-application during the fellowship. \emph{(iv) Network:} integration into the host's European and industrial collaborations, presentation at the target venues below, and engagement with civil-aviation regulators through WP5. A Career Development Plan is agreed within the first six months and formally reviewed at month~18. Target destinations after the fellowship: a tenure-track position in a European security or autonomy group, or a research leadership role in the European drone or aviation-assurance industry.

\subsection{Measures to maximise outcomes and impacts: dissemination, exploitation, communication}

\textbf{Dissemination and exploitation (specialist audiences).} Peer-reviewed publication at top security and machine-learning venues --- NDSS, IEEE Symposium on Security and Privacy, USENIX Security, IEEE SaTML --- and at networking and mobile-computing venues for the swarm work (IEEE MASS, ICDCN, IEEE Transactions on Mobile Computing), all green open access with author-accepted manuscripts deposited on publication. Beyond academia, findings are carried to the audiences that can act on them: European civil-aviation regulators and standards bodies (EASA fora on U-space, SORA and AI in aviation; the EUROCAE working groups responsible for airworthiness-security and AI assurance), the drone-operator and service-provider community through industry conferences, and the emergency-medical and environmental-agency users of the mission profiles studied. Exploitation is primarily through open infrastructure --- benchmark, harness and reference policies released under permissive licences to become the community's default evaluation substrate --- with a route to advisory or contract engagements for operators seeking assurance evidence, handled through the host's technology-transfer office.

\textbf{Communication (non-specialist audiences).} A plain-language project page and short explainer videos on why a drone carrying a defibrillator must be able to prove it will arrive. A live public demonstration of certified UAV flight under simulated interference at the \emph{Science4All} European Researchers' Night in Padova, staged with the host group, targeting school pupils and the general public; associated school-visit material on how autonomous systems are made trustworthy. Regular short-form updates through the host's channels, and open-access datasets and figures usable by science journalists. Research data are managed under a Data Management Plan following FAIR principles, with provenance tags on every released number.

\subsection{Magnitude and importance of the contribution to expected impacts}

\textbf{Scientific.} CERTIFLIGHT establishes cross-layer certified autonomy assurance as a research direction, and supplies a reusable methodology --- unify perturbation sources into one budget, compose the resulting certificate with the layer above, use it as a hard constraint on a learned controller --- that transfers directly to autonomous ground vehicles, robotic inspection and other cyber-physical systems. Expected: 4--6 peer-reviewed papers, an open benchmark and platform, and released constrained-MARL reference policies.

\textbf{Societal.} The societal case rests on three concrete civil applications, chosen because in each the mission failure mode is a safety failure, not an inconvenience. \emph{Emergency medical delivery:} drone-delivered defibrillators have arrived ahead of the ambulance in a majority of dispatches in Swedish trials, with a median lead of several minutes --- and survival in out-of-hospital cardiac arrest falls by roughly ten percent per minute of delay. A guaranteed completion floor under interference is what allows such a service to be certified for routine, unattended operation. \emph{Environmental protection:} wildfire ignition detection and flood and air-quality monitoring depend on persistent coverage; the WP4 contribution is precisely the ability to keep coverage when individual aircraft degrade. \emph{Urban monitoring and infrastructure inspection:} traffic and utility inspection over populated areas is exactly where ground-risk-driven assurance requirements bite hardest. Across all three, the widespread GNSS interference now documented over European airspace makes certified degradation behaviour a public-safety requirement rather than an academic nicety.

\textbf{Economic and regulatory.} Assurance is the binding constraint on scaling beyond-visual-line-of-sight civil drone operations, and therefore on a European commercial drone market already worth billions of euro annually and projected to roughly double by 2030. By producing evidence artefacts aligned to SORA assurance levels, EU AI Act Article~15 and DO-326A, CERTIFLIGHT gives European operators and manufacturers a concrete instrument for demonstrating robustness and cybersecurity of learned autonomy --- supporting the EU Drone Strategy 2.0, U-space deployment, and AI Act implementation, and strengthening European technological sovereignty in a market where assurance capability is a competitive differentiator.

% ==========================================================
\section{Quality and Efficiency of the Implementation}

\subsection{Work plan, work packages, deliverables, milestones and risk assessment}

\textbf{Work-package structure.} Five work packages. WP1 runs the full 24 months and carries management, training, dissemination and communication. WP2--WP4 are the scientific line, deliberately sequenced so that each supplies the next: the certificate (WP2) is the constraint consumed by the swarm controller (WP4), and the platform (WP3) is the substrate on which both are evaluated. WP5 converts results into assurance evidence.

\begin{center}
\begin{tabularx}{\linewidth}{@{}p{0.9cm}p{4.3cm}p{1.5cm}Y@{}}
\toprule
\rowcolor{LightBlue}
\textbf{WP} & \textbf{Title} & \textbf{Months} & \textbf{Objectives and main tasks} \\
\midrule
WP1 & Management, training, dissemination and communication & M1--M24 & Career Development Plan (M6, reviewed M18); data management plan; risk monitoring; open-access publication; \emph{Science4All} public demonstration; ERC-StG concept. \\
WP2 & Unified perturbation budget and composed certificates & M1--M14 & T2.1 threat and disturbance model incl.\ faults and gusts; T2.2 residual-budget lemma and interface calibration; T2.3 composition theorem and soundness; T2.4 complementary certificate suite. \\
WP3 & Cross-layer benchmark and open platform & M2--M18 & T3.1 two-layer schema and provenance flags; T3.2 ingestion adapters; T3.3 detector-operating-curve harness; T3.4 pairing pipeline; T3.5 public releases. \\
WP4 & Certified-constrained multi-agent RL for swarm re-allocation & M9--M22 & T4.1 CMDP formulation with certified radii as constraints; T4.2 Lagrangian baselines; T4.3 certified safety layer and shielded fallback; T4.4 feasibility and coverage-loss analysis; T4.5 simulation across civil mission geometries. \\
WP5 & Assurance mapping and worked certification case & M14--M24 & T5.1 mapping to SORA assurance levels, EU AI Act Art.~15, DO-326A/ED-202A; T5.2 worked assurance case for one civil mission; T5.3 regulator and industry engagement. \\
\bottomrule
\end{tabularx}
\end{center}

\textbf{Gantt overview} (shaded = active; $\blacklozenge$ = milestone).

\begin{center}
\footnotesize
\begin{tabular}{@{}l|cccccccccccc@{}}
\toprule
\textbf{WP / quarter} & Q1 & Q2 & Q3 & Q4 & Q5 & Q6 & Q7 & Q8 & & & & \\
\midrule
WP1 Mgmt/training/dissem. & \cellcolor{LightGold}$\blacklozenge$ & \cellcolor{LightGold} & \cellcolor{LightGold} & \cellcolor{LightGold}$\blacklozenge$ & \cellcolor{LightGold} & \cellcolor{LightGold} & \cellcolor{LightGold} & \cellcolor{LightGold}$\blacklozenge$ & & & & \\
WP2 Budget \& certificates & \cellcolor{LightBlue} & \cellcolor{LightBlue}$\blacklozenge$ & \cellcolor{LightBlue} & \cellcolor{LightBlue}$\blacklozenge$ & & & & & & & & \\
WP3 Benchmark \& platform & \cellcolor{LightBlue} & \cellcolor{LightBlue} & \cellcolor{LightBlue}$\blacklozenge$ & \cellcolor{LightBlue} & \cellcolor{LightBlue} & \cellcolor{LightBlue}$\blacklozenge$ & & & & & & \\
WP4 Constrained swarm RL & & & \cellcolor{LightBlue} & \cellcolor{LightBlue} & \cellcolor{LightBlue}$\blacklozenge$ & \cellcolor{LightBlue} & \cellcolor{LightBlue} & \cellcolor{LightBlue}$\blacklozenge$ & & & & \\
WP5 Assurance mapping & & & & & \cellcolor{LightBlue} & \cellcolor{LightBlue} & \cellcolor{LightBlue} & \cellcolor{LightBlue}$\blacklozenge$ & & & & \\
\bottomrule
\end{tabular}
\end{center}

\textbf{Deliverables and milestones.}
\begin{center}
\begin{tabularx}{\linewidth}{@{}p{1.3cm}p{1.1cm}Y|p{1.3cm}p{1.1cm}Y@{}}
\toprule
\rowcolor{LightBlue}
\textbf{Del.} & \textbf{Month} & \textbf{Description} & \textbf{Milest.} & \textbf{Month} & \textbf{Description} \\
\midrule
D1.1 & M6 & Career Development Plan & MS1 & M6 & CDP agreed; DMP in place \\
D1.2 & M12 & Data Management Plan (updated) & MS2 & M12 & Composition theorem proven and validated \\
D1.3 & M24 & Dissemination report; ERC-StG concept & MS3 & M18 & Benchmark v1 publicly released; CDP review \\
D2.1 & M8 & Interface calibration report & MS4 & M21 & Constrained-MARL policy meets coverage target under jamming \\
D2.2 & M14 & Composition theorem and certificate suite & MS5 & M24 & Assurance case published; platform v2 released \\
D3.1 & M6 & Schema and ingestion adapters & & & \\
D3.2 & M18 & Open benchmark and harness (v1) & & & \\
D4.1 & M16 & CMDP formulation and baselines & & & \\
D4.2 & M22 & Certified-constrained swarm policy and analysis & & & \\
D5.1 & M24 & Assurance white paper and worked case & & & \\
\bottomrule
\end{tabularx}
\end{center}

\textbf{Risk assessment and contingency.}
\begin{center}
\begin{tabularx}{\linewidth}{@{}p{4.9cm}p{1.5cm}Y@{}}
\toprule
\rowcolor{LightBlue}
\textbf{Risk} & \textbf{Likelihood / impact} & \textbf{Mitigation and contingency} \\
\midrule
Unified budget intractable as a single certificate & Med / Med & Fall back to per-source certificates combined by a union bound; the composition and the swarm layer are unaffected. \\
Interface calibration too loose to certify a useful operating window & Med / Med & The conservative kinematic bound is already sound; the tighter empirical mapping is an enhancement, not a dependency. \\
Certified-constrained MARL does not scale to target swarm size & Med / High & Staged guarantee strength: report the scalability ceiling explicitly, fall back to shielded MADDPG with empirical constraint satisfaction plus post-hoc certification of the learned policy. \\
Scarcity of same-platform measured pairings & Med / Low & Provenance flags make evidence strength explicit; physics-model and simulation pairings sustain coverage. \\
Third-party dataset licensing restricts redistribution & Low / Low & Release adapters and unified labels, fetching source data from origin under its own licence. \\
Host dependency; fellow relocation and onboarding & Low / Med & Collaboration already active with a joint manuscript in submission; host relocation support; remote start possible. \\
\bottomrule
\end{tabularx}
\end{center}

\textbf{Management and data.} The fellow manages the project day to day under the supervisor's oversight, with monthly progress reviews, a quarterly risk register review, and version-controlled research artefacts. A Data Management Plan (D1.2) governs FAIR data handling, licensing, provenance and long-term deposition. Research integrity, ethics (including dual-use assessment of UAV autonomy research and data-protection review of any imagery), and research-security requirements are handled under the host's institutional procedures and reported in Part~B2.

\subsection{Quality and capacity of the host institution and hosting arrangements}
The University of Padua is one of Europe's oldest and most research-intensive universities; the SPRITZ group is an internationally leading security and privacy laboratory with a substantial record of European project participation and of hosting funded researchers, and with the directly relevant combination of UAV network security, adversarial machine learning, and algorithms expertise. The co-supervisor's work on UAV networks, surveillance and GNSS-compromised swarm protocols is the direct counterpart to the fellow's certification expertise, and the two lines are already producing joint output.
The Host is to supply: laboratory and computing infrastructure including GPU resources and UAV testbed access; administrative, grant-office and relocation support; integration into the doctoral school and its training catalogue; the formal supervision and career-development arrangements; and the group's record of previously hosted fellows and their career outcomes.

\subsection{Eligibility}
The fellow satisfies the MSCA-PF conditions: the doctoral thesis was successfully defended in June 2026 with formal award in December 2026 (the call admits candidates who have successfully defended but not yet formally been awarded a doctoral degree at the call deadline); the fellow is well within the eight-year post-doctoral experience limit; and the mobility rule is satisfied, the fellow not having resided or carried out his main activity in Italy for more than 12 months in the 36 months immediately before the deadline. 

\vspace{0.4em}
{\footnotesize\itshape Prepared by Roger Nick Anaedevha for review with Prof.\ Mauro Conti and Dr.\ Federico Cor\`o. Not for submission yet, until you as the host confirms your sections, budget figures and you verify the entire eligibility conditions against the official 2026 call documents on the EU Funding and Tenders Portal. Links are attached to the email}

\end{document}
